\section{Open regions and their physical cuts}%
\label{sec:peps_open_regions}

A finite region of a PEPS contracts to an open tensor whose virtual legs are
the bonds leaving the region. This section introduces the open tensor, cuts
the physical vector along the boundary of a region, and defines the canonical
regions of regular tensors, with and without group elements inserted on their
bonds.

\begin{definition}[The open tensor of a finite region]%
    \label{def:peps_open_region_contraction}
    \lean{TNLean.PEPS.RegionIncidentConfig, TNLean.PEPS.regionIncidentBoundaryLabel,
        TNLean.PEPS.isRegionIncidentEdge_of_regionVertex, TNLean.PEPS.regionIncidentWeight,
        TNLean.PEPS.openRegionWeight, TNLean.PEPS.openRegionMap,
        TNLean.PEPS.openRegionMap_apply, TNLean.PEPS.openRegionMap_single,
        TNLean.PEPS.openRegionCoefficientMatrix, TNLean.PEPS.openRegionMap_apply_eq_mulVec,
        TNLean.PEPS.toMatrix_openRegionMap,
        TNLean.PEPS.regionExteriorMultiplicity}
    \leanok
    For a region $R$ in a finite PEPS graph, let $T_R(\mu;\sigma)$ be
    the contraction of its site tensors with physical configuration
    $\sigma$, fixing the crossing-bond labels $\mu$ and summing only
    the internal-bond labels. It defines the linear map from the
    crossing virtual indices to the physical system of $R$.
    No edge wholly outside $R$ enters this contraction. This is the
    finite-region tensor used when blocking and cutting the PEPS in
    \cite[Section~6, proof of Theorem~6.9]{Schuch2010PEPS}.
\end{definition}

\begin{lemma}[Exact contraction across a physical cut]%
    \label{thm:peps_open_region_cut}
    \lean{TNLean.PEPS.regionBlockedWeight_eq_exterior_mul_openRegionWeight,
        TNLean.PEPS.regionExteriorMultiplicity_eq_prod,
        TNLean.PEPS.regionInteriorBondProd_eq_exterior_mul_exterior,
        TNLean.PEPS.regionExteriorMultiplicity_ne_zero,
        TNLean.PEPS.stateCoeff_eq_openRegionComplement,
        TNLean.PEPS.stateCoeff_eq_openRegionComplement_of_regularBonds,
        TNLean.PEPS.stateCoeff_rectangle_eq_openRegionComplement_of_regularBonds}
    \leanok
    \uses{def:peps_open_region_contraction}
    Suppose every bond space has positive dimension. The coefficient
    of the PEPS at physical configurations $\sigma$ on $R$ and $\tau$
    on its complement is exactly
    \begin{align}
        \psi(\sigma,\tau)=\sum_\mu T_R(\mu;\sigma)T_{R^c}(\mu;\tau).
        \label{eq:peps_open_region_cut}
    \end{align}
    In particular this holds for regular bonds of any finite group.
    If the region coefficient is instead defined by summing global
    edge assignments, it equals $T_R$ multiplied by the product of
    the dimensions of all edges wholly outside $R$.

    This is an exact finite-region contraction identity. Identifying
    its cut maps with isometries on the regular invariant boundary
    remains a separate step in \cite[Theorem~6.9]{Schuch2010PEPS}.
\end{lemma}
\begin{proof}\leanok
    Splitting each edge assignment into its incident and exterior parts separates
    the exterior labels, which do not occur in any tensor of $R$ and therefore sum
    to the stated multiplicity. Applying the existing coefficient
    identity for the globally indexed region tensors on the two sides multiplies
    their exterior multiplicities to the product of all
    non-crossing bond dimensions; cancelling this nonzero product
    gives (\ref{eq:peps_open_region_cut}). Regular bond dimensions are $|G|>0$.
\end{proof}

\begin{definition}[Regular coordinates of the physical cut]%
    \label{def:peps_regular_physical_cut}
    \lean{TNLean.PEPS.groupBondTensor,
        TNLean.PEPS.regularLabelTranslation,
        TNLean.PEPS.regularRegionBoundaryConfigEquiv,
        TNLean.PEPS.regularOpenRegionMatrix,
        TNLean.PEPS.regularOpenComplementMatrix,
        TNLean.PEPS.regularPhysicalCutMatrix}
    \leanok
    \uses{def:peps_open_region_contraction}
    Express every bond in a group basis, and enumerate the $b=n+1$
    crossing bonds. Before the virtual reduction in
    \cite[Proof of Theorem~6.9]{Schuch2010PEPS}, let $A$ and $B$ be the open tensors of the
    region and its complement, with their columns indexed by $G^b$ in
    the same crossing-bond order, and let $K$ be the matrix of physical
    PEPS coefficients on these two systems.
\end{definition}

\begin{lemma}[Invariant regular boundary vectors]%
    \label{thm:peps_regular_boundary_invariants}
    \lean{TNLean.PEPS.regularBoundaryRepresentation_apply,
        TNLean.PEPS.regularBoundaryRepresentation,
        TNLean.PEPS.regularBoundaryRelativeEquiv,
        TNLean.PEPS.regularBoundaryRelativeEquiv_smul,
        TNLean.PEPS.regularBoundaryInvariantsEquiv,
        TNLean.PEPS.finrank_regularBoundaryInvariants_succ,
        TNLean.PEPS.finrank_regularBoundaryInvariants,
        TNLean.PEPS.regularBoundaryOrbitBasis,
        TNLean.PEPS.regularBoundaryOrbitBasis_apply,
        TNLean.PEPS.regularBoundaryOrbitBasis_eq_sum,
        TNLean.PEPS.regularBoundaryOrbitBasis_dotProduct}
    \leanok
    \uses{thm:peps_regular_blocking_coordinates}
    Let $b\geq1$ and let $G$ act on $G^b$ by simultaneous left
    multiplication. Its invariant functions have dimension $|G|^{b-1}$.
    In the invariant-boundary calculation used in the proof of
    \cite[Theorem~6.9]{Schuch2010PEPS}, writing $a=(a_0,\ldots,a_{b-1})$ makes the relative coordinates
    $(a_0^{-1}a_i)_{i=1}^{b-1}$ label the orbits, and the equal-weight orbit
    vectors form an orthogonal basis, each of squared norm $|G|$.
\end{lemma}
\begin{proof}\leanok
    The bijection $a\mapsto(a_0,(a_0^{-1}a_i)_{i>0})$ changes the action to
    left multiplication of its first coordinate alone. Invariant functions
    are therefore arbitrary functions of the remaining $b-1$ labels.
    Distinct orbits are disjoint, and each contains exactly $|G|$ elements.
\end{proof}

\begin{definition}[The virtual regular boundary state]%
    \label{def:peps_regular_boundary_state}
    \lean{TNLean.PEPS.regularBoundaryProjector,
        TNLean.PEPS.regularBoundarySchmidtMatrix,
        TNLean.PEPS.regularBoundaryState,
        TNLean.PEPS.regularBoundaryDensity,
        TNLean.PEPS.regularBoundaryTranslationState}
    \leanok
    \uses{thm:peps_regular_boundary_invariants, thm:peps_group_averaging}
    Set $b=n+1$ and $r=|G|^n$, and let $P$ be the averaging projector
    on $\mathbb C[G^b]$. The unnormalized bipartite boundary vector has
    coefficients
    \begin{align}
        \Psi(a,c)=\sum_{z\in G}\prod_{i=0}^{b-1}\delta_{a_i,zc_i}
        =|G|P_{a,c}.
    \end{align}
    For the virtual vectors used
    in the proof of \cite[Theorem~6.9]{Schuch2010PEPS}, its normalized coefficient matrix is $C=P/\sqrt r$, and its reduced
    density operator is $\rho=P/r$.
\end{definition}

\begin{lemma}[Flat spectrum of the virtual regular boundary]%
    \label{thm:peps_regular_boundary_density}
    \lean{TNLean.PEPS.regularBoundaryProjector_mul_self,
        TNLean.PEPS.rank_regularBoundaryProjector,
        TNLean.PEPS.regularBoundaryProjector_apply,
        TNLean.PEPS.regularBoundaryProjector_conjTranspose,
        TNLean.PEPS.trace_regularBoundaryProjector,
        TNLean.PEPS.regularBoundaryProjector_transpose,
        TNLean.PEPS.regularBoundarySchmidtMatrix_mul_projector,
        TNLean.PEPS.regularBoundaryProjector_mul_schmidtMatrix,
        TNLean.PEPS.regularBoundarySchmidtMatrix_mul_conjTranspose,
        TNLean.PEPS.schmidtCoeffMatrix_regularBoundaryState,
        TNLean.PEPS.partialTrace_regularBoundaryState,
        TNLean.PEPS.regularBoundaryDensity_posSemidef,
        TNLean.PEPS.trace_regularBoundaryDensity,
        TNLean.PEPS.regularBoundaryState_dotProduct,
        TNLean.PEPS.rank_regularBoundaryDensity,
        TNLean.PEPS.regularBoundaryDensity_mul_self,
        TNLean.PEPS.regularBoundaryProjector_mulVec_of_mem_invariants,
        TNLean.PEPS.regularBoundaryDensity_mulVec_of_mem_invariants,
        TNLean.PEPS.regularBoundaryTranslationState_apply,
        TNLean.PEPS.regularBoundaryTranslationState_eq_sum_prod,
        TNLean.PEPS.regularBoundaryState_eq_smul_translationState}
    \leanok
    \uses{def:peps_regular_boundary_state, thm:peps_regular_boundary_invariants}
    The averaging matrix $P$ is real, symmetric, and idempotent, with
    rank and trace $r$. Its entries are $|G|^{-1}$ on pairs in the same
    orbit and zero otherwise. The vector with coefficient matrix $C$
    has norm one, and tracing out its second factor gives
    $CC^\dagger=\rho=P/r$. Hence $\rho$ is positive, has trace one and
    rank $r$, and satisfies $\rho^2=r^{-1}\rho$. On the invariant
    boundary subspace it acts as multiplication by $r^{-1}$.
\end{lemma}
\begin{proof}\leanok
    For each pair in the same orbit, exactly one translation sends one
    configuration to the other. The entry formula follows by averaging.
    Its symmetry gives $P=P^\dagger$, and averaging twice gives $P^2=P$.
    The orbit count gives the rank; its trace follows either from the
    diagonal entries or from its projection property. Now
    $CC^\dagger=P/r$, whose trace is one. The remaining assertions follow
    from $P^2=P$ and from $P$ being the identity on invariant vectors.
\end{proof}

\begin{lemma}[Boundary symmetry and projector insertion at the cut]%
    \label{thm:peps_regular_open_region_projector}
    \lean{TNLean.PEPS.regularLabelTranslation_apply,
        TNLean.PEPS.regionIncidentBoundaryLabel_translation,
        TNLean.PEPS.regionIncidentWeight_translation,
        TNLean.PEPS.openRegionWeight_translation,
        TNLean.PEPS.regularOpenRegionMatrix_translation,
        TNLean.PEPS.matrix_mul_regularBoundaryProjector_of_translation,
        TNLean.PEPS.regularOpenRegionMatrix_mul_projector,
        TNLean.PEPS.regularPhysicalCutMatrix_eq_mul_transpose,
        TNLean.PEPS.regularPhysicalCutMatrix_eq_mul_projector_mul_transpose,
        TNLean.PEPS.regularOpenRegionMatrix_mul_schmidtMatrix_mul_transpose}
    \leanok
    \uses{def:peps_regular_physical_cut, def:peps_regular_boundary_state}
    Suppose the coefficient tensor at every site is invariant under
    simultaneous left translation of all its virtual group labels.
    Then the open tensor of every region is invariant under
    simultaneous left translation of its crossing labels. Consequently
    $AP=A$ and the physical cut satisfies
    \begin{align}
        K=AB^{\mathsf T}=APB^{\mathsf T},\qquad
        ACB^{\mathsf T}=\frac{K}{\sqrt r},
        \quad C=\frac P{\sqrt r},\quad r=|G|^n.
        \label{eq:peps_open_region_projector_cut}
    \end{align}
    In the group bases of regular and dual regular legs, this is the
    inherited boundary symmetry of the proof of
    \cite[Theorem~6.9]{Schuch2010PEPS}. It identifies the contraction
    with evaluation on the virtual boundary state, up to the scalar
    $\sqrt r$ in (\ref{eq:peps_open_region_projector_cut}). It does not assert isometry of $A,B$ or normalization of
    the physical vector.
\end{lemma}
\begin{proof}\leanok
    \uses{thm:peps_open_region_cut}
    Translating every incident-edge assignment by the same group element
    permutes the internal summation and translates the crossing
    labels, while local invariance leaves the product of site coefficients
    unchanged. Averaging this symmetry gives $AP=A$. The exact cut
    identity gives $K=AB^{\mathsf T}$, and inserting $P$ and then
    $C=P/\sqrt r$ yields the remaining formulas.
\end{proof}

\begin{lemma}[Independent finite choices]%
    \label{thm:peps_finite_indicator_sum}
    \lean{Fintype.prod_sum_boole}
    \leanok
    Let $I$ be a finite set, let $J_i$ be finite sets indexed by $i\in I$,
    and let $P_i$ be a condition on $J_i$. In every commutative semiring,
    \begin{align}
        \prod_{i\in I}\sum_{j\in J_i}\mathbf1_{P_i(j)}
        =\sum_{f\in\prod_{i\in I}J_i}
        \mathbf1_{\{P_i(f_i)\text{ for every }i\in I\}}.
    \end{align}
    Thus multiplying the separate numbers of admissible choices counts
    the admissible simultaneous choices.
\end{lemma}

\begin{proof}\leanok
    Distributing the finite product over the finite sums reduces each summand to a product of
    zero-one indicators, which is the indicator of the conjunction of their conditions.
\end{proof}

\begin{lemma}[Physical maps commute with the open-region contraction]%
    \label{thm:peps_region_physical_maps}
    \lean{TNLean.PEPS.regionPhysicalProductMatrix_mul,
        TNLean.PEPS.regionPhysicalProductMatrix,
        TNLean.PEPS.regionPhysicalProductMatrix_one,
        TNLean.PEPS.regionPhysicalProductMatrix_conjTranspose,
        TNLean.PEPS.regionPhysicalProductMatrix_isIsometry,
        TNLean.PEPS.regionPhysicalMap,
        TNLean.PEPS.regionPhysicalMap_apply,
        TNLean.PEPS.regionPhysicalMap_comp,
        TNLean.PEPS.regionPhysicalMap_openRegionWeight}
    \leanok
    \uses{def:peps_open_region_contraction}
    For a region $R$, let $F_v$ be a physical matrix at each vertex, allowing
    different input and output dimensions at different vertices. Their product
    matrix is
    \begin{align}
        F_R(\tau,\sigma)=\prod_{v\in R}F_v(\tau_v,\sigma_v).
    \end{align}
    Then $(FL)_R=F_RL_R$, $(F^\dagger)_R=F_R^\dagger$, and identity local maps
    give the identity on the physical space of the region. In particular,
    local isometries give an isometry $F_R$.

    For the physical operation on a block in
    \cite[\S6.1, accessible virtual coordinates]{Schuch2010PEPS}, applying $F_R$ to the open-region tensor is the contraction of the
    transformed site coefficients:
    \begin{align}
        (F_R T_R)(\tau,\mu)=
        \sum_{\eta:\eta|_{\partial R}=\mu}
        \prod_{v\in R}\sum_s F_v(\tau_v,s)A_v(s,\eta|_v).
    \end{align}
\end{lemma}
\begin{proof}\leanok
    Expanding the matrix products and interchanging the finite sums over physical
    and virtual labels factors the sum over independent physical labels at
    each vertex. Composition follows from the same finite distributivity;
    complex conjugation gives the adjoint identity. Applying these identities
    to $F_v^\dagger F_v=\mathbb 1$ proves the isometry assertion.
\end{proof}

\begin{definition}[Local regular Gram kernel]%
    \label{def:peps_regular_site_gram}
    \lean{TNLean.PEPS.regularLegRepresentation,
        TNLean.PEPS.regularLegRepresentation_fin,
        TNLean.PEPS.regularLegRepresentation_apply,
        TNLean.PEPS.regularLegRepresentation_unitary,
        TNLean.PEPS.regularLegProjector,
        TNLean.PEPS.regularLegProjector_fin,
        TNLean.PEPS.regularLegProjector_apply,
        TNLean.PEPS.regularSiteMap,
        TNLean.PEPS.toMatrix_regularSiteMap,
        TNLean.PEPS.tensorSiteGram}
    \leanok
    Let $E_v$ be the set of virtual legs incident to a site $v$.
    In the group basis, its coefficient map $A_v$ has entries
    $A_v(s,\eta)=a_v(\eta;s)$, for $\eta\in G^{E_v}$.
    Simultaneous left translation acts on functions by
    $(U_gf)(\eta)=f(g^{-1}\eta)$, and its invariant projector has entries
    \begin{align}
        P_v(\eta,\theta)=\frac1{|G|}\sum_{g\in G}
        \mathbf 1_{\{\eta=g\theta\}}.
    \end{align}
    The local Gram kernel is
    $\sum_s\overline{a_v(\eta;s)}a_v(\theta;s)$.
\end{definition}

\begin{lemma}[Gram formula at an isometric regular site]%
    \label{thm:peps_regular_site_gram}
    \lean{TNLean.PEPS.IsGIsometric.exists_regularSiteGram,
        TNLean.PEPS.IsGInjective.regularSiteMap_translation}
    \leanok
    \uses{def:peps_regular_site_gram}
    If $A_v$ is $G$-isometric, there is $c_v>0$ such that
    \begin{align}
        \sum_s\overline{a_v(\eta;s)}a_v(\theta;s)
        =\frac{c_v}{|G|}\sum_{g\in G}
        \mathbf 1_{\{\eta=g\theta\}}.
        \label{eq:peps_regular_site_gram}
    \end{align}
    Its coefficients are invariant under simultaneous left translation;
    this invariance already follows from $G$-injectivity.
    This is the group-basis form of
    \cite[Definition~6.1]{Schuch2010PEPS}, with the positive scalar
    convention of \cite{gap:rmp_peps_quantum_double_g_isometry}.
\end{lemma}
\begin{proof}\leanok
    The defining adjoint identity is $A_v^\dagger A_v=c_vP_v$.
    Evaluating both sides in the group basis gives
    (\ref{eq:peps_regular_site_gram}).
    The identity $A_vU_g=A_v$ gives coefficient invariance.
\end{proof}

\begin{definition}[The canonical regular-projector region]%
    \label{def:peps_regular_projector_open_region}
    \lean{TNLean.PEPS.RegionHalfEdgeConfig,
        TNLean.PEPS.regularProjectorOpenRegionMatrix,
        TNLean.PEPS.regularSiteMap_regularLegProjector,
        TNLean.PEPS.isGIsometric_regularLegProjector,
        TNLean.PEPS.regularProjectorOpenRegionMatrix_apply,
        TNLean.PEPS.regionIncidentBoundaryLabel_regularGroup_iff}
    \leanok
    \uses{def:peps_open_region_contraction, def:peps_regular_site_gram,
        thm:peps_finite_indicator_sum}
    Let $R$ be a finite region of the original graph, $E_R$ its incident
    edges, and $\partial R$ its crossing edges. Write $E_v$ for the
    incident edges at $v$ and put $\kappa=|G|$. The canonical physical
    configuration is $\alpha=(\alpha_v)_{v\in R}$, with
    $\alpha_v\in G^{E_v}$. The local canonical site map is the regular
    averaging projector $P_v$; it is $G$-isometric with factor one.
    Define the canonical open-region coefficient matrix by
    \begin{align}
        B_R(\alpha,\theta)
         & =\sum_{\substack{\eta\in G^{E_R}                    \\
                \eta|_{\partial R}=\theta}}
        \prod_{v\in R}P_v(\alpha_v,\eta|_{E_v})\notag          \\
         & =\kappa^{-|R|}\sum_{\eta\in G^{E_R}}\sum_{q\in G^R}
        \mathbf1_{\{\eta|_{\partial R}=\theta\}}
        \mathbf1_{\{\alpha_v=q_v\eta|_{E_v}\ (v\in R)\}}.
        \label{eq:peps_projector_region}
    \end{align}
    Every internal bond is summed once, and no exterior bond label
    occurs. As in the exposed virtual coordinates
    of \cite[Observation after Lemma~6.3]{Schuch2010PEPS}, encoding group labels in finite coordinates preserves
    the boundary condition.
\end{definition}

\begin{lemma}[Accessibility of a regular region]%
    \label{thm:peps_regular_open_region_accessibility}
    \lean{TNLean.PEPS.IsGInjective.exists_regularProjectorCoefficients,
        TNLean.PEPS.IsGIsometric.exists_regularProjectorCoefficients,
        TNLean.PEPS.IsGInjective.regularSiteMap_projector_coefficients,
        TNLean.PEPS.regionPhysicalMap_eq_regularProjectorOpenRegionMatrix,
        TNLean.PEPS.regionPhysicalProductMatrix_mul_regularOpenRegionMatrix,
        TNLean.PEPS.regionPhysicalMap_regularProjectorOpenRegionMatrix,
        TNLean.PEPS.exists_regularProjectorOpenRegionMatrix}
    \leanok
    \uses{def:peps_regular_projector_open_region, def:peps_g_injective}
    Suppose the local regular site maps $A_v$ are $G$-isometric,
    with positive factors $c_v$. Let $T_R(\theta)$ be the
    contracted physical column with crossing labels $\theta$,
    and write $B_R(\theta)$ for the column in
    Definition~\ref{def:peps_regular_projector_open_region}. With
    $L_v=c_v^{-1}A_v^\dagger$, one has
    \begin{align}
        \left(\bigotimes_{v\in R}L_v\right)T_R(\theta)
         & =B_R(\theta),\notag \\
        \left(\bigotimes_{v\in R}A_v\right)B_R(\theta)
         & =T_R(\theta).
        \label{eq:peps_region_accessibility}
    \end{align}
    The reverse identity needs only the local invariance of the
    $G$-injective site maps. More generally, any local matrices
    $F_v$ satisfying $F_vA_v=P_v$ give the forward identity.
    Such matrices also exist for arbitrary regular $G$-injective sites,
    by their local left-inverse property.
    These assertions concern the contracted columns and
    do not require the physical maps to be inverse on unused
    coordinates. No region Gram identity is assumed.
\end{lemma}

\begin{proof}\leanok
    \uses{thm:peps_regular_site_gram, thm:peps_region_physical_maps}
    The local identities are
    $c_v^{-1}A_v^\dagger A_v=P_v$ and $A_vP_v=A_v$.
    The first follows from the local regular Gram kernel; the second
    follows from virtual invariance by averaging the group action.
    Applying the physical product maps inside the internal-bond
    contraction turns each local coefficient into the corresponding
    matrix product, giving~(\ref{eq:peps_region_accessibility}).
    Expanding $P_v$ and distributing the finite products gives
    the translation formula~(\ref{eq:peps_projector_region}).
\end{proof}

\begin{definition}[Regular bonds with inserted group operators]%
    \label{def:peps_regular_twisted_region}
    \lean{TNLean.PEPS.regularTwistedLabels,
        TNLean.PEPS.regularTwistedSite,
        TNLean.PEPS.regularTwistedOpenRegionMatrix}
    \leanok
    \uses{def:peps_regular_physical_cut}
    Order the endpoints of each bond $e$, writing them as $s(e)<t(e)$.
    For an assignment $u_e\in G$, the inserted regular matrix makes
    the label at $t(e)$ equal to $u_e$ times the label at $s(e)$.
    Equivalently, replacing the coefficient at vertex $v$ by
    \begin{align}
        A^u_{v,s}(\eta)=A_{v,s}(\eta^u),\qquad
        \eta^u_e=
        \begin{cases}u_e\eta_e, & v=t(e), \\ \eta_e,&v=s(e),
        \end{cases}
    \end{align}
    and contracting the internal bonds with a fixed crossing configuration
    gives the open map $T_R^u$.
\end{definition}

\begin{definition}[The canonical projector region with bond insertions]%
    \label{def:peps_regular_projector_twisted_region}
    \lean{TNLean.PEPS.regularProjectorTwistedRegionMatrix,
        TNLean.PEPS.regularProjectorTwistedRegionMatrix_one,
        TNLean.PEPS.regularProjectorTwistedRegionMatrix_apply}
    \leanok
    \uses{def:peps_regular_twisted_region,
        def:peps_regular_projector_open_region, thm:peps_finite_indicator_sum}
    For arbitrary inserted group elements $u_e$, including those on
    crossing bonds, define
    \begin{align}
        B_R^u(\alpha,\theta)
         & =\sum_{\substack{\eta\in G^{E_R}                    \\
                \eta|_{\partial R}=\theta}}
        \prod_{v\in R}P_v(\alpha_v,\eta^u|_{E_v})              \\
         & =\kappa^{-|R|}\sum_{\eta\in G^{E_R}}\sum_{q\in G^R}
        \mathbf1_{\{\eta|_{\partial R}=\theta\}}
        \mathbf1_{\{\alpha_v=q_v\eta^u|_{E_v}\ (v\in R)\}}.
        \label{eq:peps_twisted_projector_region}
    \end{align}
    Here the half-edge modification $\eta^u$ is that of
    Definition~\ref{def:peps_regular_twisted_region}. Identity insertions
    recover $B_R$ as the finite contraction of the local
    averaging-projector sites with the inserted regular operators.
\end{definition}

\begin{lemma}[Accessibility with arbitrary regular bond insertions]%
    \label{thm:peps_regular_twisted_region_accessibility}
    \lean{TNLean.PEPS.regionPhysicalMap_eq_regularProjectorTwistedRegionMatrix,
        TNLean.PEPS.regionPhysicalMap_regularProjectorTwistedRegionMatrix,
        TNLean.PEPS.exists_regularProjectorTwistedRegionMatrix_of_isGInjective,
        TNLean.PEPS.exists_regularProjectorTwistedRegionMatrix}
    \leanok
    \uses{def:peps_regular_projector_twisted_region}
    Let the original regular site maps $A_v$ be $G$-injective, and
    choose local matrices $F_v$ satisfying $F_vA_v=P_v$. The same
    matrices work for every region and every assignment of inserted
    group elements:
    \begin{align}
        \left(\bigotimes_{v\in R}F_v\right)T_R^u(\theta)=B_R^u(\theta),
        \qquad
        \left(\bigotimes_{v\in R}A_v\right)B_R^u(\theta)=T_R^u(\theta).
    \end{align}
    If the sites are $G$-isometric with factors $c_v>0$, one may take
    $F_v=c_v^{-1}A_v^\dagger$. No connectedness or global Gram formula
    is required. These are the accessibility identities of
    \cite[Observation after Lemma~6.3]{Schuch2010PEPS}, extended
    algebraically to arbitrary bond insertions.
\end{lemma}

\begin{proof}\leanok
    \uses{thm:peps_regular_open_region_accessibility,
        thm:peps_region_physical_maps}
    Applying the physical product map inside the contraction preserves
    the local identity $F_vA_v=P_v$ at the modified
    virtual label $\eta^u|_{E_v}$, proving the forward equality.
    In the reverse direction, $A_vP_v=A_v$ gives
    $A_vF_vA_v=A_v$ at the modified label, and expanding the local
    averages yields the finite-sum formula
    (\ref{eq:peps_twisted_projector_region}) for $B_R^u$.
\end{proof}

\begin{lemma}[The torus closure as a twisted graph contraction]%
    \label{thm:peps_regular_torus_closure_cut_identification}
    \lean{TNLean.PEPS.torusClosureEdgeAssignment_right,
        TNLean.PEPS.torusIncidentSite,
        TNLean.PEPS.torusClosureEdgeAssignment,
        TNLean.PEPS.torusClosureEdgeAssignment_up,
        TNLean.PEPS.torusRightEdge_head_of_wrap,
        TNLean.PEPS.torusUpEdge_head_of_wrap,
        TNLean.PEPS.torusRegularBondConfigEquiv,
        TNLean.PEPS.torusClosureBondConfigEquiv,
        TNLean.PEPS.torusClosureBondConfigEquiv_right,
        TNLean.PEPS.torusClosureBondConfigEquiv_up,
        TNLean.PEPS.torusClosureTwistedLabels_top,
        TNLean.PEPS.torusClosureTwistedLabels_right,
        TNLean.PEPS.torusClosureTwistedLabels_down,
        TNLean.PEPS.torusClosureTwistedLabels_left,
        TNLean.PEPS.torusGClosure_eq_stateCoeff_twisted,
        TNLean.PEPS.torusGClosure_assembleRegion_eq_regularPhysicalCutMatrix,
        TNLean.PEPS.torusGClosure_cut_eq_regularPhysicalCutState}
    \leanok
    \uses{def:peps_torus_g_closure, def:peps_regular_twisted_region}
    On a torus of width and height at least three, the four incident
    graph edges are exactly the top, right, down and left virtual legs.
    The closure vector $\psi_{g,h}$ equals the contraction of
    the graph tensor with ordered-head assignment
    \begin{align}
        u_e=
        \begin{cases}
            h^{-1}, & e\text{ is a horizontal wrap bond}, \\
            g,      & e\text{ is a vertical wrap bond},   \\
            1,      & \text{otherwise}.
        \end{cases}
    \end{align}
    Its coefficient vector across any region cut is the twisted
    physical cut state. No commutativity or isometry assumption is
    required for these contraction identities. The circumference
    restriction ensures that the simple graph distinguishes all four
    incident virtual legs \cite{gap:rmp_peps_quantum_double_g_isometry}.
\end{lemma}

\begin{proof}\leanok
    Enumerating the graph bonds by horizontal and vertical torus bonds,
    at a horizontal wrap bond the ordered head is the native tail;
    replacing its summed label by $h$ times the native tail label makes the inserted $h^{-1}$ give the original native tail, while
    the other endpoint gives the native head. At a vertical wrap bond,
    the ordered head is the native head, so the assignment is $g$.
    All four site labels then coincide with those of the native closure
    contraction under a bijective change of labels, and reading the resulting
    identity in the two sets of physical cut coordinates gives the cut statements.
\end{proof}
